\section{Warriors and Actualized Men}

\begin{quotex}
This is a guest post by Rory, who had been working out a personal program. He proposes that it, in whole or in part, may be useful for a new Fourth Order. I have registered the domain The Bernard Option\footnote{\url{https://www.bernardoption.com/}}, although it has not yet been configured. If there is enough interest, I will move similar discussions to the new site. If you would like to contribute your own ideas, please let me know. I also hope someone would be interested in assuming the duties of a webmaster. 
\flright{Administrator}
\end{quotex}

A strange bit of synchronicity occurred with your Templar post. For a while, I have been sketching and contemplating a sort of program that I feel the past and perhaps future Templar order would have had in the process of training warriors and actualized men. I've drawn up seven disciplines based on various sources such as \textbf{Marcus Aurelius}, \textbf{Geoffrey de Charny}, and \textbf{Baldassare Castiglione}`s \emph{Book of the Courtier} and perhaps more than a few personal idiosyncrasies.

Briefly:

\paragraph{Asceticism/the Greek Training}


Marcus Aurelius makes mention of the Greek training which were ascetic disciplines put forward by different philosophical schools. This would be a system of ascetic disciplines covering physical training, diet, sex, and way of life.

\paragraph{War}


This would be a combination of individual combat training of striking, grappling, weapons, and firearms. Alongside would be a study of tactics, strategy, and leadership

\paragraph{Survival}


A grounding in fieldcraft\footnote{\url{https://en.wikipedia.org/wiki/Fieldcraft}} and emergency response/medical skills as well a general ability to be at home in nature if the need arose.

\paragraph{Profession}


The world being what it is this would be the money generating discipline in the program. While the original Templars were involved in finance I'm partial to a kind of gentleman farmer archetype\footnote{\url{http://distributistreview.com/is-distributism-agrarianism/}}. I'm still developing whether this would be a trade or some other form of business.

\paragraph{Social Dynamics}


This would be a skill in negotiating and diplomacy. You catch more flies with honey than vinegar and external considering allows a level of understanding that would resolve frustrations and aid in self-reflection.

\paragraph{Liberal Arts}


For thousands of years the trivium and quadrivium\footnote{\url{http://cassiodorusquodlibeta.blogspot.com/2016/05/why-higher-education-is-anything-but.html}} were considered the basis of an educated gentleman. A grounding in these arts would open one's mind to the order pervading the world.

\paragraph{Contributing to the Hermetic stream}


This would be the spiritual sun around which the disciplines would orbit. Spiritual devotion and practice bridging the churches of Peter and John. Ideally, the disciplines would draw from this and serve as vehicles to a greater awareness of the spirit and build individuals as pathfinders of Tradition and carry the Hermetic torch.

Of course, any one of these disciplines could entail decades of study so a jack of all trades approach would be the best one could hope for with a brief tour of this mortal coil. But also believe all the disciplines would coalesce and overlap in myriad ways.

I'd been tossing these ideas about and your post triggered a new perspective. I had also been thinking of ways to incorporate the minor arcana to structure the training. Of course, I hope to avoid the temptation to build but rather plant good seeds in the soul and ``add style to one's character''. Another obstacle is to avoid another modern bourgeoisie system of self-improvement. I'd like to believe my discernment of spirits is somewhat adequate.


\flrightit{Posted on 2017-07-24 by Rory}

\begin{center}* * *\end{center}

\begin{footnotesize}
\begin{sffamily}
\texttt{yourusernameistakeng on 2017-07-26 at 21:50 said:}

I really enjoyed this post. I had no idea about distributist theory, and will absolutely keep researching this. I think that economy wise, agrarianism is where it is at nowadays. Not only farming, but finding ways to farm in unusual places such as forests and other areas. From my training in ``higher education'' which dealt with biotechnology and plant genes extensively, I gather there has been very low interest in natural farming innovation. Instead we want to splice a gene and spend a couple 100,000 on technology that does this.

\hfill

\texttt{Aleksandar on 2017-07-27 at 04:48 said:}

I would like to contribute more but I'm not so proficient in english and also view myself as a learner, not a teacher. I can always provide my commentaries which are not very useful and smart and always come out as some kind of train of thought ruminations. What i'm trying to say is -- these recent posts really aligned with my own interests and recent thoughts. I feel very interested and wish I could contribute more in any way possible.

Living in the rural area for most of my life so far, in the middle of nowhere, helped me learn a thing or two about farming. One of my life goals is actually owning a farm, which is contradictory to most people here who's dream is escaping the rural and blending into the urban.

Even now I often visit my relative's farm and try to absorb as much information as possible. I always thought it's weird on my side to, in one hand, have interest in ``intellectual'' (hate a modern connotation of this word) aspects of life -- studying history, esotericism, theology and philosophy -- which now I simply synthesize and categorize as Tradition and, on the other, wanting to realize myself as a farmer. The need of being self-sufficient and largely independent through one's own endeavor in this day and age is almost looked upon as something primitive, culturally obsolete. It's easier to just give in.

I also think that any actualized man, especially a warrior type, should get into a real fistfight at one point in time, earlier the better. I speak of nothing sinister of course, just a good old ``pub brawl'' type of situation. I'm not suggesting starting an incident directly, but there are ways and life brings many opportunities for such harmless exercise in fundamental principles of masculinity.

Martial arts are very useful, but if not possible, just practicing simple punching at home is enough. Robert E. Howard is a good archetype here. An amateur boxer, lifting weights, somewhat traditional, intellectual type. Lifting is always helpful, but of course remaining within the limits of every day functionality. If you can't scratch your back anymore you're overdoing it. I'm lifting at home when I have time and it's more of an exercise of willpower than strength. Building a core strength and classic aesthetics is more important than bodybuilding a statue for looking good on a beach.

Greek training combined with ascetic disciplines is ideal. Roman legionnaires also had a good training. I love those willpower exercises like showering in cold water, fasting, walking long distances, enduring all night without sleep and so on.

I think maybe having some kind of summer camps would be great. Studying esoteric matter, camping, working out, reading poetry and having general discussions. Also creating a network of individuals around the world who would serve as some kind of vigilants/knights, but in what exact way, i'm not yet so sure.

\hfill

\texttt{yourusernameistakeng on 2017-07-28 at 06:28 said:}

I would also like to suggest that maybe creating a distributist sort of technology for currency amongst members of this new ``Bernard Option'' might be useful. I'm wary of that for the longterm. Everyone might need to pool the knowledge they have depending on their professions to really kickstart this thing.

\hfill

\texttt{Logres on 2017-08-02 at 23:20 said:}

I would suggest the quarter staff as an excellent Western style martial arts weapon. Not only is it seemingly an innocuous staff or pole for carrying things on it, but there are several documentary accounts historically of it being superior (in the hands of some one reasonably well trained) over a sword (wielded by a similar opponent).
\url{https://www.linkedin.com/pulse/brief-history-quarterstaff-frank-docherty}

``The Duke smiled scornfully and beckoned a third man to join the original two. Peeke and the rapier men warily traversed each other, all the while thrusting and warding, till finally Peeke gambled on an all out attack. His first blow a left one of

\hfill

\texttt{Logres on 2017-08-02 at 23:21 said:}

his adversaries dead and his subsequent blows left the other two injured and disarmed. No doubt they also left the Spanish seriously questioning the wisdom of their invasion plans. Peeke's feat so impressed his Spanish captors that they released him and granted him safe conduct to England.'' Hector Paulous Mair has a Renaissance treatise on all types of Western weapons martial arts combat, and you can find men testing all of them out in various re-stagings of the sword plays described by Mair. Quarterstaffs and longbows aren't exactly high powered ``threats'' in the modern world, but then again, people who train with them are just re-enactors, and not really a concern. In any event, it is the development of courage in the face of danger and daring in the face of evil that is the main quarry -- the rest will come.

\end{sffamily}
\end{footnotesize}

